% ID: SEL-D0008-BB-E04
% ESTADO: cotejada. Segunda pasada visual completa y borrador compilado revisados.
% Fuente: Selectividad/matematicas_aplicadas/ORD_SUPLENTE_Modelo-4-Enunciados (Mat-Ap-2021).pdf
% SHA-256: 74b067b81be20aa0262a0cc58795f3b25488aa9c749d825b7b33f656dd1ae089
% Página: 1; bloque B; ejercicio original 4.

Una fábrica estima que sus costes de producción, expresados en miles de
euros, vienen dados por la función $f(x)=x^2-6x+10$, donde $x$ es la cantidad
semanal a producir expresada en miles de kilogramos.
\begin{enumerate}[label=\alph*)]
  \item \textbf{(1 punto)} \textquestiondown{}Cuál debe ser la producción
  semanal para que el coste sea mínimo? \textquestiondown{}Cuál es dicho coste?
  \item \textbf{(1.5 puntos)} Calcule la recta tangente a la función de costes
  en el punto de abscisa $x=4$. Represente gráficamente la función de costes
  y la recta tangente hallada.
\end{enumerate}
